% sec:dg-geodesic — The Geodesic Equation
\section{The Geodesic Equation}
\label{sec:dg-geodesic}

A freely falling particle --- one subject to no forces other than gravity ---
follows the straightest possible path through curved spacetime.  In flat
spacetime, ``straightest'' means ``straight line'': a particle at rest stays at
rest, and a moving particle continues at constant velocity.  On a curved
manifold, straight lines don't exist, but the concept of ``straightest''
survives as a curve whose tangent vector doesn't turn relative to the geometry
itself.  The geodesic equation makes this precise, and every photon trajectory
computed by the ray tracer is a solution of it.
\physnote{``Freely falling'' in GR means subject only to gravity.  Gravity
is not a force --- it is the curvature of spacetime.  A particle in free fall
follows a geodesic; a particle that feels a non-gravitational force does not.}

\paragraph{Parallel transport and autoparallels.}

Consider a curve $x^\mu(\lambda)$ parametrised by $\lambda$, with tangent
vector $\dot{x}^\mu = \d x^\mu / \d\lambda$.  A vector field $V^\mu$ defined
along this curve is \emph{parallel-transported} if its covariant rate of change
in the direction of the tangent vanishes:
%
\begin{equation}
  \dot{x}^\alpha\,\cd{\alpha} V^\mu = 0 \,.
  \label{eq:parallel-transport}
\end{equation}
%
Expanding the covariant derivative using \cref{eq:covariant-deriv-vector}
gives the component form
%
\begin{equation}
  \frac{\d V^\mu}{\d\lambda}
  + \christoffel{\mu}{\alpha}{\beta}\,\dot{x}^\alpha\, V^\beta = 0 \,.
  \label{eq:parallel-transport-components}
\end{equation}
%
The first term is the ordinary rate of change of $V^\mu$ along the curve; the
second corrects for the rotation of the coordinate basis, exactly as the
covariant derivative was designed to do in \cref{sec:dg-covariant}.

A geodesic is a curve that parallel-transports its own tangent vector --- it
is as straight as the geometry allows.  Setting $V^\mu = \dot{x}^\mu$ in
\cref{eq:parallel-transport-components} yields the \emph{geodesic equation}:
%
\begin{equation}
  \ddot{x}^\mu + \christoffel{\mu}{\alpha}{\beta}\,
  \dot{x}^\alpha\,\dot{x}^\beta = 0 \,,
  \label{eq:geodesic}
\end{equation}
%
where $\ddot{x}^\mu = \d^2 x^\mu / \d\lambda^2$.  This is a system of four
coupled second-order ODEs for the coordinate functions $x^\mu(\lambda)$.  The
Christoffel symbols encode a position-dependent ``force'' that bends the
trajectory, but it is not a physical force --- it is the geometry itself,
determined by the metric derivatives in \cref{eq:christoffel-formula}, that
shapes freely falling worldlines.
\xrefnote{The geodesic equation is the curved-spacetime generalisation of
the photon trajectory equation from \cref{sec:sr-causal}.}

\paragraph{The variational principle.}

The same equation arises from a different starting point: extremising the
action along a curve.  Define the Lagrangian
%
\begin{equation}
  L = \tfrac{1}{2}\,\metric\,\dot{x}^\mu\,\dot{x}^\nu \,,
  \label{eq:geodesic-lagrangian}
\end{equation}
%
so that $2L$ is the squared norm of the tangent vector.  The factor of
$\tfrac{1}{2}$ simplifies the algebra.  Demanding that the action
$S = \int L\,\d\lambda$ be stationary under variations with fixed endpoints
yields the Euler--Lagrange equations
%
\begin{equation}
  \frac{\d}{\d\lambda}\frac{\partial L}{\partial \dot{x}^\mu}
  - \frac{\partial L}{\partial x^\mu} = 0 \,.
  \label{eq:euler-lagrange}
\end{equation}
%
The canonical momentum is
$\partial L / \partial\dot{x}^\mu = \metric\,\dot{x}^\nu$, and the
position derivative is
$\partial L / \partial x^\mu = \tfrac{1}{2}\,(\pd{\mu}\,
g_{\alpha\beta})\,\dot{x}^\alpha\,\dot{x}^\beta$.  Substituting into
\cref{eq:euler-lagrange} and expanding
$\tfrac{\d}{\d\lambda}(\metric\,\dot{x}^\nu)$ with the product rule yields
three metric-derivative terms.  Since $\dot{x}^\alpha\,\dot{x}^\beta$ is
symmetric, these combine into the three-term Christoffel combination
$\pd{\alpha}\,g_{\mu\beta} + \pd{\beta}\,g_{\mu\alpha}
- \pd{\mu}\,g_{\alpha\beta}$ from \cref{eq:christoffel-formula}; contracting
with the inverse metric recovers \cref{eq:geodesic}.

The two derivations --- parallel transport and the variational principle ---
yield the same equation.  This is a theorem, not a coincidence: for a
metric-compatible, torsion-free connection (the Levi-Civita connection of
\cref{sec:dg-covariant}), autoparallels and extremals coincide
\cite{Carroll2004}.
\histnote{The equivalence between variational and parallel-transport geodesics
depends on the torsion-free condition.  In Einstein--Cartan theory, which
permits torsion, the two notions diverge.}

\paragraph{Affine parametrisation and the null--timelike split.}

\Cref{eq:geodesic} holds only for \emph{affine} parameters --- those related
to any other affine parameter by a linear transformation
$\lambda \to a\lambda + b$ with $a \neq 0$.  The Lagrangian $L$ is constant
along any affinely parametrised geodesic, which gives the causal
classification:
%
\begin{equation}
  2L = \metric\,\dot{x}^\mu\,\dot{x}^\nu =
  \begin{cases}
    -1 & \text{timelike ($\lambda = \tau$, proper time)} \\
    \phantom{-}0 & \text{null (photons)} \\
    +1 & \text{spacelike}
  \end{cases}
  \label{eq:geodesic-constraint}
\end{equation}
%
For timelike geodesics, the natural affine parameter is proper time $\tau$, and
the tangent vector $\dot{x}^\mu = \fourvel$ is the four-velocity with
$\metric\,u^\mu\,u^\nu = -1$.  For null geodesics, proper time vanishes along
the worldline, so the affine parameter $\lambda$ has no direct physical
interpretation --- it can be rescaled by any positive constant without changing
the trajectory.  The ray tracer exploits this freedom by setting the photon
energy $E = 1$, as in \cref{sec:sr-causal}.
\warnnote{The constraint \cref{eq:geodesic-constraint} is a first integral,
not an additional equation.  If the initial tangent vector satisfies the
constraint, the geodesic equation preserves it automatically.}

\paragraph{The Hamiltonian formulation.}

The geodesic equation is a second-order system: four equations for four
unknowns $x^\mu(\lambda)$.  Numerical integrators prefer first-order systems,
since standard adaptive methods (Runge--Kutta, symplectic integrators) are
designed for that form.  The Legendre transform of the Lagrangian defines the
conjugate momentum
%
\begin{equation}
  p_\mu = \frac{\partial L}{\partial \dot{x}^\mu}
        = \metric\,\dot{x}^\nu \,,
  \label{eq:conjugate-momentum}
\end{equation}
%
which is the index-lowered tangent vector --- the same covariant four-momentum
introduced in \cref{sec:sr-causal}.  The Hamiltonian is $\mathcal{H} =
p_\mu\,\dot{x}^\mu - L$; substituting
$\dot{x}^\mu = \inversemetric\,p_\nu$ and simplifying gives the
\emph{geodesic super-Hamiltonian}:
%
\begin{equation}
  \mathcal{H} = \tfrac{1}{2}\,\inversemetric\,p_\mu\,p_\nu \,.
  \label{eq:geodesic-hamiltonian}
\end{equation}
%
Hamilton's equations then provide eight first-order ODEs:
%
\begin{align}
  \dot{x}^\mu &= \frac{\partial\mathcal{H}}{\partial p_\mu}
               = \inversemetric\,p_\nu \,,
  \label{eq:hamilton-xdot} \\
  \dot{p}_\mu &= -\frac{\partial\mathcal{H}}{\partial x^\mu}
               = -\tfrac{1}{2}\,(\pd{\mu}\,g^{\alpha\beta})\,
                 p_\alpha\,p_\beta \,.
  \label{eq:hamilton-pdot}
\end{align}
%
The first equation is the algebraic inversion of
\cref{eq:conjugate-momentum} using the non-degenerate inverse metric.  The
second shows that $\dot{p}_\mu$ depends on partial derivatives of the metric
components, not on Christoffel symbols directly.
\coderef{Geodesic.Hamiltonian}{hamiltonianRHS}

An equivalent form of the momentum equation follows directly from the
Euler--Lagrange equation: $\dot{p}_\mu = \partial L / \partial x^\mu$ gives
%
\begin{equation}
  \dot{p}_\mu = \tfrac{1}{2}\,(\pd{\mu}\,g_{\alpha\beta})\,
    \dot{x}^\alpha\,\dot{x}^\beta \,,
  \label{eq:hamilton-pdot-covariant}
\end{equation}
%
which involves derivatives of the \emph{covariant} metric rather than the
inverse.  The sign change relative to \cref{eq:hamilton-pdot} arises from the
identity $\pd{\mu}\,g^{\alpha\beta} =
-g^{\alpha\gamma}\,g^{\beta\delta}\,\pd{\mu}\,g_{\gamma\delta}$ (obtained by
differentiating $g^{\alpha\gamma}\,g_{\gamma\beta} = \delta^\alpha{}_\beta$).
The codebase uses this form: the automatic differentiation strategy from
\cref{sec:dg-covariant} differentiates the covariant metric function directly,
and \cref{eq:hamilton-pdot-covariant} consumes those derivatives without ever
requiring the inverse metric to be differentiated.  Hamilton's equations thus
need only the metric, its inverse, and first derivatives of the metric.
\implnote{The function \texttt{hamiltonianRHS} computes $\dot{x}^\mu$ via
\cref{eq:hamilton-xdot}, then evaluates \cref{eq:hamilton-pdot-covariant}
using AD-computed metric derivatives.}

Metric compatibility ($\cd{\alpha}\,\metric = 0$) guarantees that the
Lagrangian ($\dot{x}^\mu$) and Hamiltonian ($p_\mu$) formulations are fully
interchangeable --- the result promised in \cref{sec:dg-covariant}.

\paragraph{Constraint preservation.}

Since $\mathcal{H}$ depends on $x^\mu$ and $p_\mu$ but not on $\lambda$
explicitly, the chain rule applied to Hamilton's equations gives
$\dot{\mathcal{H}} = \partial\mathcal{H}/\partial\lambda = 0$: the
Hamiltonian is conserved along any solution.  Comparing with the normalisation
conditions in \cref{eq:geodesic-constraint}:
%
\begin{equation}
  \mathcal{H} =
  \begin{cases}
    -\tfrac{1}{2} & \text{timelike} \\
    \phantom{-}0 & \text{null}
  \end{cases}
  \label{eq:hamiltonian-constraint}
\end{equation}
%
For null geodesics, $\mathcal{H} = 0$ is preserved exactly by the continuous
equations; any drift during numerical integration is entirely due to
discretisation error.  The codebase monitors $\abs{\mathcal{H}}$ at every
step --- for a well-tuned adaptive integrator, this quantity remains below
$10^{-12}$ throughout a typical ray trace, providing a built-in accuracy
diagnostic at no extra computational cost.
\xrefnote{The full numerical integration strategy --- adaptive step control,
Killing-vector conservation, and termination conditions --- is developed in
\cref{ch:geodesics}.}

\paragraph{A worked example: great circles on $S^2$.}

The geodesic equation on the two-sphere, with the Christoffel symbols of
\cref{eq:sphere-christoffel}, gives two component equations:
%
\begin{align}
  \ddot{\theta} - \sin\theta\cos\theta\;\dot{\varphi}^2 &= 0 \,,
  \label{eq:sphere-geodesic-theta} \\
  \ddot{\varphi} + 2\cot\theta\;\dot{\theta}\,\dot{\varphi} &= 0 \,.
  \label{eq:sphere-geodesic-phi}
\end{align}
%
Consider the equator: $\theta(\lambda) = \pi/2$ for all $\lambda$, with
$\varphi(\lambda) = \omega\lambda$ for constant angular speed $\omega$.  Then
$\dot{\theta} = \ddot{\theta} = 0$ and $\ddot{\varphi} = 0$.
\Cref{eq:sphere-geodesic-theta} requires
$-\sin(\pi/2)\cos(\pi/2)\,\omega^2 = 0$, which holds since
$\cos(\pi/2) = 0$.  \Cref{eq:sphere-geodesic-phi} reduces to $0 = 0$.  The
equator is a geodesic --- a great circle traversed at constant angular speed.
By contrast, a latitude circle at $\theta_0 \neq \pi/2$ fails
\cref{eq:sphere-geodesic-theta} because $\cos\theta_0 \neq 0$ --- the curve
must constantly accelerate poleward to maintain its latitude, confirming that
only great circles are geodesics on $S^2$.
\physnote{Every great circle on $S^2$ is a geodesic.  A rotation of
coordinates maps any great circle to the equator, so verifying the equator
suffices by symmetry.}

\begin{figure}[htbp]
  \centering
  % sphere-geodesics.tex — Geodesics (great circles) on the two-sphere S²
% Inputable TikZ fragment for fig:sphere-geodesics
% Requires: tikz with calc, arrows.meta (loaded by ehbook.cls)
%
% Projection: orthographic, viewing elevation = tilt above equatorial plane.
%
% For a point at colatitude θ, azimuth φ on the sphere:
%   screen_x = R·sinθ·cosφ
%   screen_y = R·sinθ·sinφ·sin(tilt) + R·cosθ·cos(tilt)
%   depth    = R·sinθ·sinφ·cos(tilt) - R·cosθ·sin(tilt)  [> 0 = far side]
%
% A latitude circle at colatitude θ projects to an ellipse centred at
% (0, R·cosθ·cos(tilt)) with semi-axes (R·sinθ, R·sinθ·sin(tilt)).
% The far-side arc (behind the sphere) spans ellipse angles limbT to
% 180°−limbT, where sinφ_limb = cotθ·tan(tilt).
%
\begin{tikzpicture}[
    >={Stealth[length=5pt]},
    font=\footnotesize,
  ]

  % ── Parameters ──────────────────────────────────────────────────────
  \def\R{1.8}                              % sphere radius (cm)
  \pgfmathsetmacro{\tilt}{30}              % elevation above equatorial plane (°)
  \pgfmathsetmacro{\eqb}{\R*sin(\tilt)}   % equator semi-minor axis

  % ── Sphere outline and shading ──────────────────────────────────────
  \shade[ball color=EHMarginGray!8, draw=EHMarginGray, thin]
    (0,0) circle (\R);

  % ── Great circle 1: equator ─────────────────────────────────────────
  % Far side (behind sphere)
  \draw[EHAccretionGold, thick, opacity=0.3]
    ({-\R},0) arc (180:360:{\R} and {\eqb});
  % Near side
  \draw[EHAccretionGold, thick]
    ({\R},0) arc (0:180:{\R} and {\eqb});

  % ── Great circle 2: tilted ~35° clockwise in the picture plane ──────
  \begin{scope}[rotate=35]
    \draw[EHAccretionGold, thick, opacity=0.3]
      ({-\R},0) arc (180:360:{\R} and {\R*0.15});
    \draw[EHAccretionGold, thick]
      ({\R},0) arc (0:180:{\R} and {\R*0.15});
  \end{scope}

  % ── Great circle 3: tilted ~40° anticlockwise ──────────────────────
  \begin{scope}[rotate=-40]
    \draw[EHAccretionGold, thick, opacity=0.3]
      ({-\R},0) arc (180:360:{\R} and {\R*0.18});
    \draw[EHAccretionGold, thick]
      ({\R},0) arc (0:180:{\R} and {\R*0.18});
  \end{scope}

  % ── Latitude line at θ = π/4 (non-geodesic) ────────────────────────
  %
  % Correct orthographic projection of a latitude circle at colatitude θ:
  %   centre:     (0, R·cosθ·cos(tilt))
  %   semi-major: R·sinθ               (horizontal, no foreshortening)
  %   semi-minor: R·sinθ·sin(tilt)     (vertical, foreshortened)
  %
  % The far-side arc (behind the limb) spans ellipse parameter
  % t ∈ [limbT, 180°−limbT], where limbT = arcsin(cotθ · tan(tilt)).
  % At the limb-crossing points the projected curve sits exactly on
  % the sphere outline (distance R from origin).
  %
  \pgfmathsetmacro{\latR}{\R*sin(45)}
  \pgfmathsetmacro{\latY}{\R*cos(45)*cos(\tilt)}      % centre y
  \pgfmathsetmacro{\latb}{\latR*sin(\tilt)}            % semi-minor
  \pgfmathsetmacro{\limbT}{asin(cos(45)/sin(45)*tan(\tilt))}  % ≈ 35.3°

  % Far-side arc (behind the sphere): limbT to 180−limbT
  \draw[EHMarginGray, densely dashed, thick, opacity=0.3]
    ({\latR*cos(\limbT)}, {\latY + \latb*sin(\limbT)})
    arc ({\limbT}:{180-\limbT}:{\latR} and {\latb});

  % Near-side arc (visible): 180−limbT through 360 back to limbT
  \draw[EHMarginGray, densely dashed, thick]
    ({-\latR*cos(\limbT)}, {\latY + \latb*sin(\limbT)})
    arc ({180-\limbT}:{360+\limbT}:{\latR} and {\latb});

  % ── Tangent vector on the equator (geodesic) ───────────────────────
  % Point at ~65° on the near-side equator arc
  \pgfmathsetmacro{\eqPx}{\R*cos(65)}
  \pgfmathsetmacro{\eqPy}{\eqb*sin(65)}
  % Tangent direction: derivative of ellipse parametrisation at t=65°
  \pgfmathsetmacro{\eqTx}{-sin(65)}
  \pgfmathsetmacro{\eqTy}{\eqb/\R*cos(65)}
  \pgfmathsetmacro{\eqTlen}{sqrt(\eqTx*\eqTx + \eqTy*\eqTy)}
  \pgfmathsetmacro{\eqTxn}{0.7*\eqTx/\eqTlen}
  \pgfmathsetmacro{\eqTyn}{0.7*\eqTy/\eqTlen}
  \draw[->, EHJetBlue, very thick]
    (\eqPx, \eqPy) -- ++(\eqTxn, \eqTyn)
    node[above left, inner sep=1pt, font=\scriptsize, text=EHJetBlue]
    {$\dot{x}^\mu$};

  % ── Tangent vector on the latitude line (non-geodesic) ─────────────
  % Point at ~330° (on the visible near-side arc, right of centre)
  \pgfmathsetmacro{\ltAng}{330}
  \pgfmathsetmacro{\ltPx}{\latR*cos(\ltAng)}
  \pgfmathsetmacro{\ltPy}{\latY + \latb*sin(\ltAng)}
  % Tangent direction along the latitude ellipse
  \pgfmathsetmacro{\ltTx}{-sin(\ltAng)}
  \pgfmathsetmacro{\ltTy}{\latb/\latR*cos(\ltAng)}
  \pgfmathsetmacro{\ltTlen}{sqrt(\ltTx*\ltTx + \ltTy*\ltTy)}
  \pgfmathsetmacro{\ltTxn}{0.65*\ltTx/\ltTlen}
  \pgfmathsetmacro{\ltTyn}{0.65*\ltTy/\ltTlen}
  \draw[->, EHJetBlue, very thick]
    (\ltPx, \ltPy) -- ++(\ltTxn, \ltTyn)
    node[above right, inner sep=1pt, font=\scriptsize, text=EHJetBlue]
    {$\dot{x}^\mu$};

  % ── Labels ─────────────────────────────────────────────────────────

  % Equator label — outside the sphere to the lower-right
  \node[font=\scriptsize, text=EHAccretionGold!80!black, anchor=west]
    at ({\R + 0.12}, {-\eqb*0.4}) {equator};

  % Latitude label — near the top of the dashed arc, above vectors
  \node[font=\scriptsize, text=EHMarginGray, anchor=south]
    at ({0.4}, {\latY + \latb + 0.08})
    {$\theta = \pi/4$};

  % Sphere label
  \node[font=\small] at (0, {-\R - 0.35}) {$S^2$};

\end{tikzpicture}

  \caption[Geodesics on the two-sphere]{Geodesics on the two-sphere are great circles --- curves whose
    centres coincide with the centre of the sphere.  Three great circles
    (solid gold) are shown, including the equator.  The dashed grey curve
    is a latitude line at $\theta = \pi/4$; it is not a geodesic because
    a parallel-transported tangent vector rotates away from the curve's
    tangent direction.  All spheres share the same geodesic structure
    regardless of radius, since the factor $R^2$ cancels from the
    geodesic equation.}
  \label{fig:sphere-geodesics}
\end{figure}

The geodesic equation, in either its second-order form (\cref{eq:geodesic}) or
the Hamiltonian system (\cref{eq:hamilton-xdot,eq:hamilton-pdot}), is the
equation of motion that the ray tracer integrates at every step.  The
Hamiltonian form is preferred computationally: it yields a first-order system
whose conserved Hamiltonian provides a built-in error diagnostic, and any
coordinate absent from the metric produces a conserved component of $p_\mu$
directly from \cref{eq:hamilton-pdot}.  The remaining geometric ingredient is a
measure of how the gravitational field curves nearby geodesics apart --- the
subject of the next section.
